\documentclass[11pt]{article}
\usepackage[margin=1in]{geometry}
\usepackage{amsmath,amssymb}
\usepackage{enumitem}
\usepackage{hyperref}
\usepackage{mathtools}

\title{2D Derivation Target for the 2022 Magnetic-Field-Inspired Navigation Paper}
\author{}
\date{2026-06-18}

\begin{document}
\maketitle

\section*{Purpose}

This note fixes a concrete 2D target for the clean Python implementation.
The 2022 paper is written in a general 3D vector-field form, but the current
benchmark platform is planar. The goal of this document is to define a precise
2D reduction target for:

\begin{itemize}[leftmargin=*]
  \item the robot motion direction \(l_a\),
  \item the obstacle surface normal \(n_o\),
  \item the artificial surface current \(l_o\),
  \item the perpendicular / opposite current \(l_o^\perp\),
  \item the boundary-following field \(F_b\),
  \item the collision-avoidance field \(F_a\).
\end{itemize}

This document is not yet a proof that the current code matches the paper.
It is the derivation target that the code should match.

\section*{Paper Context}

The 2022 paper states:

\begin{itemize}[leftmargin=*]
  \item the robot has a local sensor of radius \(r_s\),
  \item the robot velocity direction is \(l_a\),
  \item the closest obstacle point is described by vector \(r_o\),
  \item the obstacle normal is taken opposite to \(r_o\),
  \item the artificial current \(l_o\) is the projection of \(l_a\) onto the local obstacle surface,
  \item \(F_b\) is induced by \(l_o\),
  \item \(F_a\) is induced by \(l_o^\perp\),
  \item both \(F_b\) and \(F_a\) should be orthogonal to robot motion and therefore not change robot speed.
\end{itemize}

\section*{2D Setup}

We work in \(\mathbb{R}^2\). Let:

\begin{equation}
p \in \mathbb{R}^2
\end{equation}

denote the robot position,

\begin{equation}
v = \dot{p} \in \mathbb{R}^2
\end{equation}

denote the robot velocity, and

\begin{equation}
r_o \in \mathbb{R}^2
\end{equation}

denote the vector from the robot to the currently selected closest obstacle
point.

The robot speed is

\begin{equation}
\lVert v \rVert
\end{equation}

and the robot motion direction is

\begin{equation}
l_a =
\begin{cases}
\dfrac{v}{\lVert v \rVert}, & \lVert v \rVert > 0 \\
0, & \lVert v \rVert = 0
\end{cases}
\end{equation}

\section*{Local Surface Normal and Tangent}

Following the paper assumption, the obstacle surface normal at the currently
sensed closest point is taken opposite to \(r_o\):

\begin{equation}
n_o = -\frac{r_o}{\lVert r_o \rVert}
\end{equation}

whenever \(\lVert r_o \rVert > 0\).

The orthogonal projector onto the local obstacle tangent line is therefore

\begin{equation}
P_t = I - n_o n_o^\top
\end{equation}

\section*{Artificial Surface Current \(l_o\)}

The paper defines \(l_o\) as the projection of the robot motion direction
\(l_a\) onto the local obstacle surface.

In 2D, the natural reduction is:

\begin{equation}
\tilde{l}_o = P_t l_a = l_a - (l_a^\top n_o)n_o
\end{equation}

If \(\lVert \tilde{l}_o \rVert \ge \epsilon\), the unit artificial current is

\begin{equation}
l_o = \frac{\tilde{l}_o}{\lVert \tilde{l}_o \rVert}
\end{equation}

If \(\lVert \tilde{l}_o \rVert < \epsilon\), the paper says to use the unit
vector of the previously proposed current. Therefore the 2D target rule is:

\begin{equation}
l_o =
\begin{cases}
\dfrac{\tilde{l}_o}{\lVert \tilde{l}_o \rVert}, & \lVert \tilde{l}_o \rVert \ge \epsilon \\
l_{o,\text{prev}}, & \lVert \tilde{l}_o \rVert < \epsilon
\end{cases}
\end{equation}

\section*{2D Cross-Product Reduction}

To reduce the paper's 3D cross-product expressions into a planar form, define
the 90-degree rotation matrix

\begin{equation}
J =
\begin{bmatrix}
0 & -1 \\
1 & 0
\end{bmatrix}
\end{equation}

and define the scalar 2D cross product

\begin{equation}
a \times_2 b = a_x b_y - a_y b_x
\end{equation}

for \(a,b \in \mathbb{R}^2\).

If \(a,b \in \mathbb{R}^2\) are embedded into \(\mathbb{R}^3\) with zero
third component, then:

\begin{equation}
a \times b = (0,0,a \times_2 b)
\end{equation}

and

\begin{equation}
a \times (0,0,\beta) = -\beta J a
\end{equation}

These identities are the basis of the exact 2D reduction target below.

\section*{2D Target for \(F_b\)}

The paper boundary-following field follows the charged-particle style magnetic
construction:

\begin{equation}
B_b = \frac{l_o \times v}{\lVert r_o \rVert}
\end{equation}

\begin{equation}
F_b = c \, l_a \times B_b
\end{equation}

Reducing this to 2D gives the target:

\begin{equation}
B_b^{(z)} = \frac{l_o \times_2 v}{\lVert r_o \rVert}
\end{equation}

\begin{equation}
F_b^{2D} = -c \, B_b^{(z)} \, J l_a
\end{equation}

or equivalently,

\begin{equation}
F_b^{2D}
=
-c
\frac{l_o \times_2 v}{\lVert r_o \rVert}
J l_a
\end{equation}

\textbf{Important property}

Because \(J l_a\) is orthogonal to \(l_a\),

\begin{equation}
(F_b^{2D})^\top l_a = 0
\end{equation}

which is the 2D target expression for the paper statement that \(F_b\) does not
change robot speed.

\section*{2D Target for \(F_a\)}

After checking the old ROS \texttt{field=2} branch more carefully, the best 2D
target for the current clean implementation is not a second magnetic
cross-product law. The stronger implementation evidence is:

\begin{itemize}[leftmargin=*]
  \item \(F_b\) is the magnetic boundary-following term,
  \item \(F_a\) is a short-range radial collision term,
  \item \(F_a\) is projected to remain orthogonal to the motion direction so it
  does not change speed.
\end{itemize}

Define the raw short-range repulsion

\begin{equation}
\tilde{F}_a
=
-c_\perp
\left(
\frac{1}{\lVert r_o \rVert} - \frac{1}{r_{la}}
\right)
\frac{r_o}{\lVert r_o \rVert}
\end{equation}

for \(\lVert r_o \rVert \le r_{la}\). The working 2D target is then the
orthogonal projection of this repulsion onto the subspace perpendicular to
\(l_a\):

\begin{equation}
F_a^{2D}
=
\left(I - l_a l_a^\top \right)\tilde{F}_a
\end{equation}

or equivalently,

\begin{equation}
F_a^{2D}
=
\tilde{F}_a - \left(\tilde{F}_a^\top l_a\right)l_a
\end{equation}

\textbf{Important property}

By construction,

\begin{equation}
(F_a^{2D})^\top l_a = 0
\end{equation}

which is the 2D target expression for the paper statement that \(F_a\) does not
change robot speed.

\section*{Thresholded Activation Target}

In implementation, the paper uses:

\begin{itemize}[leftmargin=*]
  \item \(F_b\) only when \(\lVert r_o \rVert \le r_l\)
  \item \(F_a\) only when \(\lVert r_o \rVert \le r_{la}\)
\end{itemize}

So the 2D paper target is:

\begin{equation}
F_b^{\text{active}} =
\begin{cases}
F_b^{2D}, & \lVert r_o \rVert \le r_l \\
0, & \text{otherwise}
\end{cases}
\end{equation}

\begin{equation}
F_a^{\text{active}} =
\begin{cases}
F_a^{2D}, & \lVert r_o \rVert \le r_{la} \\
0, & \text{otherwise}
\end{cases}
\end{equation}

\section*{Local Sensing Target}

The paper assumes:

\begin{itemize}[leftmargin=*]
  \item only currently visible local sensor points are available
  \item the closest point is chosen from those readings
  \item if several close readings fall within threshold \(\delta r\), averaging
  is used for the special non-unique closest-point case
\end{itemize}

Therefore the 2D implementation target is:

\begin{enumerate}[leftmargin=*]
  \item collect only sensor-visible obstacle points within radius \(r_s\),
  \item select the minimum-distance sensed vector,
  \item form the averaged vector over readings with distance less than
  \(d_{\min} + \delta r\),
  \item use the averaged vector only when the paper Section 4.4 condition is satisfied.
\end{enumerate}

\section*{What This Means For The Next Implementation Step}

The next implementation step should treat the following as the exact 2D target:

\begin{itemize}[leftmargin=*]
  \item \(n_o = -r_o/\lVert r_o \rVert\)
  \item \(l_o = \mathrm{normalize}\!\left((I-n_on_o^\top)l_a\right)\), with the
  \(\epsilon\)-memory rule
  \item \(F_b^{2D} = -c \dfrac{l_o \times_2 v}{\lVert r_o \rVert} J l_a\)
  \item \(F_a^{2D} = \left(I-l_a l_a^\top\right)\tilde{F}_a\), where
  \(\tilde{F}_a = -c_\perp\left(\dfrac{1}{\lVert r_o \rVert} -
  \dfrac{1}{r_{la}}\right)\dfrac{r_o}{\lVert r_o \rVert}\)
\end{itemize}

This target may still need revision after checking more carefully against the
paper equations. However, it now matches the old ROS implementation evidence
for the short-range collision term much better than the rejected magnetic
\(l_o^\perp\)-based target.

\end{document}
